% 選択式問題（AMC 形式）の雛形
%
% - AMC 本体がなくても、このファイルをビルドすれば問題用紙の PDF（確認用）が作れます。
% - マークシートの印刷・スキャン・自動採点は AMC 本体で行います（docs/SETUP_AMC.md）。
\documentclass[a4paper,11pt]{ltjsarticle}
\usepackage{lfs-fonts}
\usepackage{amsmath,amssymb}
\usepackage[nowatermark]{automultiplechoice}
\AMCrandomseed{1234}   % 乱数の種。問題・選択肢の並びを固定する
% 設問の見出しを「問1」形式にする（既定は英語の "Question 1"）
\def\AMCbeginQuestion#1#2{\par\noindent{\bfseries 問#1}\quad #2}

\begin{document}

%%% 問題の定義 %%%
\element{basic}{
  \begin{question}{Q1}
    $2^3 + 3^2$ の値は？
    \begin{choices}
      \wrongchoice{$12$}
      \correctchoice{$17$}
      \wrongchoice{$15$}
      \wrongchoice{$36$}
    \end{choices}
  \end{question}
}

\element{basic}{
  \begin{question}{Q2}
    $f(x) = x^2$ のとき $f'(3)$ の値は？
    \begin{choices}
      \wrongchoice{$3$}
      \correctchoice{$6$}
      \wrongchoice{$9$}
      \wrongchoice{$2$}
    \end{choices}
  \end{question}
}

%%% 用紙のレイアウト %%%
\onecopy{1}{
  \noindent{\large\bfseries サンプル科目　第01回　確認問題}\hfill
  氏名\,\underline{\hspace{45mm}}\par\medskip
  \AMCcodeGridInt[h]{student}{8}   % 学籍番号（8桁）をマークする欄
  \medskip
  \shufflegroup{basic}
  \insertgroup{basic}
}

\end{document}
